\documentclass{article}
\usepackage{amsmath,amssymb}
\usepackage[margin=1in]{geometry}
\allowdisplaybreaks
\usepackage[UTF8]{ctex}
\begin{document}

\section*{完整的 VRPSSR 马尔可夫决策过程}

\[
\max_{\pi\in\Pi}\;
\mathbb{E}\left[\left.\sum_{k=0}^{K}
R\bigl(s_k,X_k^{\pi}(s_k)\bigr)\,\right|\,s_0\right]
\]
受每个决策规则的约束，该规则决定了在适用的每个状态下应做出何种可行的决策：
\[
X_k^{\pi}(s_k)\in\mathcal{X}(s_k),\qquad k=0,\ldots,K.
\]
可行的决策集如下所打印的块：
\begin{align}
\mathcal{X}(s_k)=\bigl\{(a,c):\quad & \nonumber\\[-1.5ex]
&a\in\{0,1\}^{|\tilde{\mathcal{C}}_k|}, \tag{1}\\
&c\in\bar{\mathcal{C}}_k\cup\tilde{\mathcal{C}}'_k
       \cup\{c_k\}\cup\{\text{depot}\}, \tag{2}\\
&c\neq\text{depot}\quad\text{if}\quad
 \bar{\mathcal{C}}_k\cup\tilde{\mathcal{C}}'_k\setminus\{c_k\}
 \neq\emptyset, \tag{3}\\
&\text{feasible routing}\ \bigr\}. \tag{4}
\end{align}

\paragraph{赋值和路由的含义。}
决策就是这样的一个配对。
\[
x=(a,c),
\]
其中二进制向量$a$有$|\tilde{\mathcal{C}}_k|$个分量。每个分量在等于一时会将请求分配给车辆，在等于零时则拒绝该请求。被拒绝的请求不会重新考虑。位置分量$c$指示车辆前往服务区域$\mathcal{C}$内的某个位置。新分配的客户集合是
\[
\tilde{\mathcal{C}}'_k
=\left\{c\in\tilde{\mathcal{C}}_k:
 a_{\tilde{\mathcal{C}}_k^{-1}(c)}=1\right\},
\]
其中，$\tilde{\mathcal{C}}_k^{-1}(c)$返回的是元素$c$在数组$\tilde{\mathcal{C}}_k$中的索引。

设定$c=c_k$是指当前位置等待基础时间单位$\bar t$。条件（3）在源代码中描述为：当$\bar{\mathcal{C}}_k$和$\tilde{\mathcal{C}}'_k$中的客户尚未得到服务时，不允许前往仓库。条件（4）要求从当前位置出发，经过所有分配的客户，再返回仓库，且所需时间小于或等于剩余时间$T-t_k$减去在当前位置等待的时间。这是一个路线存在性的要求，而不是最便宜插入测试。源代码中将这一要求表述为“可行的路线规划”，而没有提供基于弧线的路线规划公式。

\paragraph{印刷范围（3）。}
(3)中的表达式保持不变，没有在集合差之前的并集周围加括号。附带的文字解释了上述含义；印刷表达并没有明确界定$\{c_k\}$所代表的减法范围。

\section*{固定的输入、状态以及边界条件}

有一辆车辆，其工作时间为$T$天，服务区域为$\mathcal{C}$，还有一处维修站。服务请求在工作日内的任意时间出现，且出现时间在$\mathcal{C}$以内。那些提前提出请求的客户，在工作日开始之前就已经被知晓了。
\[
\mathcal{C}_{\mathrm{early}}\subseteq\mathcal{C}.
\]
该车辆必须服务于这些客户，并且必须在仓库处开始和结束服务。目标是最大化服务那些在一天中随时提交请求的客户的预期数量。

固定函数$d(\cdot,\cdot)$可以计算出$\mathcal{C}$中两个位置之间的已知旅行时间。固定数值$\bar t$是基础等待时间单位。这里没有针对特定实例的旅行时间公式、等待时间值、请求分配或数值持续时间。

在决策时刻$k$时，状态为
\[
s_k=(c_k,t_k,\bar{\mathcal{C}}_k,\tilde{\mathcal{C}}_k),
\qquad k=0,\ldots,K,
\]
与…一起
\[
c_k\in\mathcal{C},\qquad
 t_k\in[0,T],\qquad
 \bar{\mathcal{C}}_k\subseteq\mathcal{C},\qquad
 \tilde{\mathcal{C}}_k\subseteq\mathcal{C}.
\]
这里$c_k$是车辆的位置，代表客户地点或仓库；$t_k$是车辆到达$c_k$的时间，标志着决策时期$k$的开始；而$\bar{\mathcal{C}}_k$是尚未服务的分配客户集合。对于初始时期之后的时期，$\tilde{\mathcal{C}}_k$是在或早于$t_k$但晚于$t_{k-1}$时提出的、可能为空的服务请求集合，即在
\[
(t_{k-1},t_k],\qquad k=1,\ldots,K,
\]
受以下最终请求假设的约束。

初始状态是
\[
s_0=(\text{depot},0,\mathcal{C}_{\mathrm{early}},\emptyset).
\]
作为输入的可行性假设，存在一条从仓库出发，经过$\mathcal{C}_{\mathrm{early}}$中的每个客户，再回到仓库的路线，且该路线的耗时小于或等于$T$。

最终决策时刻$K$可能是一个随机变量。在这个时刻，该过程会处于一个终端状态。
\[
s_K\in
\left\{(\text{depot},t_K,\emptyset,\emptyset):t_K\in[0,T]\right\}.
\]
因此，该车辆在时间$T$时已返回仓库，并且已经为所有提前预约的客户提供了服务。来源方明确假设，最后的预约请求列表为空。
\[
\tilde{\mathcal{C}}_K=\emptyset.
\]

\section*{奖励与动态定义}

当决策$x$在状态$s_k$下被做出时，奖励金额等于被分配到的迟交订单客户的数量。
\[
R(s_k,x)=\bigl|\tilde{\mathcal{C}}'_k(s_k,x)\bigr|.
\]
这里，$\tilde{\mathcal{C}}'_k(s_k,x)$表示由$s_k$和$a$定义的集合$\tilde{\mathcal{C}}'_k$，以及$x$的赋值组件。

在$s_k$中选择$x$会将该过程推进到决策后的状态。
\[
s_k^x=(c_k,t_k,\bar{\mathcal{C}}_k^x),
\qquad
\bar{\mathcal{C}}_k^x
=\bar{\mathcal{C}}_k\cup\tilde{\mathcal{C}}'_k.
\]
被分配给客户的请求集因此被更新，以包含新分配的请求。接下来的决策前状态是……
\[
s_{k+1}=(c_{k+1},t_{k+1},\bar{\mathcal{C}}_{k+1},
\tilde{\mathcal{C}}_{k+1}).
\]
对于过渡到下一个时代的情况，$k=0,\ldots,K-1$能够区分移动和等待。

\paragraph{移动：$c\neq c_k$。}
决策时刻$k+1$在到达位置$c$时开始。
\[
c_{k+1}=c,\qquad
 t_{k+1}=t_k+d(c_k,c_{k+1}),\qquad
 \bar{\mathcal{C}}_{k+1}
 =\bar{\mathcal{C}}_k^x\setminus\{c_k\}.
\]
最后一个关系反映了车辆在之前的位置$c_k$的服务情况，而不是到达目的地后的服务情况。

\paragraph{等待中：$c=c_k$。}
该车辆停留在当前位置，决策时刻$k+1$在$\bar t$之后开始。源程序如下：
\[
t_{k+1}=t_k+\bar t,\qquad
\bar{\mathcal{C}}_k=\bar{\mathcal{C}}_k^x.
\]
\emph{打印出的不一致性。}等待段落中提到，分配给客户的集合保持不变，但实际上打印出来的却是$\bar{\mathcal{C}}_k=\bar{\mathcal{C}}_k^x$，并且两边都有下标$k$，而不是对$\bar{\mathcal{C}}_{k+1}$的更新。与决策后的定义相比，这种等式与在等待过程中添加真正新的分配情况相矛盾。因此，这里的等式被保留下来，而不需要替换其左侧的下标，也不需要提供缺失的分配集合更新信息。

在$k+1$这一时刻，可能会观察到一个新的、可能为空的客户集合$\tilde{\mathcal{C}}_{k+1}$，这些客户是在$(t_k,t_{k+1}]$期间提出的请求，其中$\tilde{\mathcal{C}}_K=\emptyset$是明确的例外情况。总体上的随机请求规律仍然保持抽象；在计算实验中使用的具体分布并不是这个完整模型的条件。

\section*{策略类别与范围}

政策是一系列决策规则的集合。
\[
\pi=(X_0^{\pi},X_1^{\pi},\ldots,X_K^{\pi}).
\]
该来源通过以下方式描述每条规则：
\[
X_k^{\pi}(s_k):s_k\mapsto\mathcal{X}(s_k),
\qquad k=0,\ldots,K,
\]
这意味着当前的状态被映射到一个可行的决策上。优化变量$\pi$的取值范围包括$\Pi$，即所有马尔可夫确定性策略的集合。目标中的期望值是以$s_0$为条件的，它考虑了随机请求以及可能随机出现的最终时刻$K$。

\paragraph{来源与选定范围。}
目标函数和策略类见 Ulmer、Goodson、Mattfeld 与 Hennig 的论文 \emph{随机请求下动态车辆路径问题的离线—在线近似动态规划}，日期为 2017 年 4 月 4 日，PDF 实际第 8 至 9 页。 输入、状态和边界定义在第6-7页；可行决策块、奖励和转移过程在第8页。 第 39 至 40 页的附录中描述了一种减少决策空间的算法，该算法不使用等待时间，而是采用有序的顾客分配方式，并且以最便宜的方式插入顾客。 该简化过程并不属于所选的完整MDP的一部分，因此不会替代条件(4)。

\end{document}
